% Wave-14 B3/20. Positive geometry of quiver gluing:
% the C\v ech chamber structure of the quiver atlas on
% $\mathrm{Stab}(X)$ as a positive-geometry / cluster-algebra datum
% on the moduli of CY$_3$ stability conditions.
%
% Companions (consulted in full before inscription):
%   notes/wave17_c1_conifold_gluing_explicit.tex (Klebanov--Witten
%     two-chart gluing, overlap cocycle = Seiberg duality).
%   notes/wave18_e5_cluster_Yplus.tex (Amiot cluster category, Iyama
%     --Yoshino cluster-tilting, Keller--Yang mutation = derived
%     equivalence, path-independence-up-to-inner-automorphism on
%     $\Yplus$).
%   notes/wave17_d4_gL_positive_geom_family.tex (lattice-polarised
%     reduced DT positive effective divisor $D_{\Phi_L}$).
%   notes/wave17_d5_preface_positive_geometry_synthesis.tex
%     (universal grammar $\Yplus(X) = \Heq(\Meff(X))$).
%   notes/wave12_a10_wall_crossing_soibelman.tex,
%   notes/wave13_a10_DT_wallcrossing_soibelman.tex (Kontsevich--
%     Soibelman wall-crossing, motivic DT, cluster transformations).
%   notes/CoHA\_to\_W\_infty\_treatise.tex (positive-half CoHA as
%     Yangian positive half $\Yplus(\ghat)$ for $\C^3$).
%   notes/platonic\_synthesis\_waves\_11\_through\_16.tex
%     (seven-face $r_{\mathrm{CY}}$ consolidation).
%
% Primary sources (all consulted at first-principles detail):
%   Fomin, S., Zelevinsky, A., "Cluster algebras I: Foundations",
%     J.\ Amer.\ Math.\ Soc.\ 15 (2002) 497--529 (arXiv:math/0104151).
%   Fomin, S., Zelevinsky, A., "Cluster algebras II: Finite type
%     classification", Invent.\ Math.\ 154 (2003) 63--121
%     (arXiv:math/0208229).
%   Fock, V., Goncharov, A., "Cluster ensembles, quantization and the
%     dilogarithm", Ann.\ Sci.\ \'ENS 42 (2009) 865--930
%     (arXiv:math/0311245).
%   Gross, M., Hacking, P., Keel, S., Kontsevich, M.,
%     "Canonical bases for cluster algebras", J.\ Amer.\ Math.\ Soc.\ 31
%     (2018) 497--608 (arXiv:1411.1394).
%   Keller, B., Yang, D., "Derived equivalences from mutations of
%     quivers with potential", Adv.\ Math.\ 226 (2011) 2118--2168
%     (arXiv:0906.0761).
%   Iyama, O., Reiten, I., "Fomin--Zelevinsky mutation and tilting
%     modules over Calabi--Yau algebras", Amer.\ J.\ Math.\ 130 (2008)
%     1087--1149 (arXiv:math/0605136).
%   Derksen, H., Weyman, J., Zelevinsky, A., "Quivers with potentials
%     and their representations I: Mutations", Selecta Math.\ 14 (2008)
%     59--119 (arXiv:0704.0649).
%   Kontsevich, M., Soibelman, Y., "Stability structures, motivic
%     Donaldson--Thomas invariants and cluster transformations"
%     (arXiv:0811.2435).
%   Kontsevich, M., Soibelman, Y., "Cohomological Hall algebra,
%     exponential Hodge structures and motivic DT invariants"
%     (arXiv:1006.2706).
%   Bridgeland, T., "Stability conditions on triangulated categories",
%     Ann.\ of Math.\ 166 (2007) 317--345 (arXiv:math/0212237).
%   Nagao, K., "Donaldson--Thomas theory and cluster algebras",
%     Duke Math.\ J.\ 162 (2013) 1313--1367 (arXiv:1002.4884).
%   Arkani-Hamed, N., Bai, Y., Lam, T., "Positive geometries and
%     canonical forms", J.\ High Energy Phys.\ 11 (2017) 039
%     (arXiv:1703.04541).
%   Arkani-Hamed, N., Trnka, J., "The amplituhedron", J.\ High Energy
%     Phys.\ 10 (2014) 030 (arXiv:1312.2007).
%   Hacking, P., Keel, S., "Mirror symmetry and cluster algebras",
%     Proc.\ ICM Rio 2018, vol.\ II, pp.\ 671--697 (arXiv:1712.05385).
%   Joyce, D., Song, Y., "A theory of generalized Donaldson--Thomas
%     invariants", Mem.\ Amer.\ Math.\ Soc.\ 217, no.\ 1020 (2012)
%     (arXiv:0810.5645).
%   Reineke, M., "Poisson automorphisms and quiver moduli",
%     J.\ Inst.\ Math.\ Jussieu 9 (2010) 653--667 (arXiv:0804.3214).
%   Szendr\H{o}i, B., "Non-commutative Donaldson--Thomas theory and
%     the conifold", Geom.\ Topol.\ 12 (2008) 1171--1202
%     (arXiv:0705.3419).
%
% Author: Raeez Lorgat. No AI attribution anywhere.

\providecommand{\ClaimStatusTheorem}{\textsc{[Theorem]}}
\providecommand{\ClaimStatusConjectured}{\textsc{[Conjectural]}}
\providecommand{\ClaimStatusDerivation}{\textsc{[First-principles]}}
\providecommand{\ClaimStatusDefinition}{\textsc{[Definition]}}
\providecommand{\ClaimStatusSynthesis}{\textsc{[Synthesis]}}

\providecommand{\K}{\mathbf{k}}
\providecommand{\Z}{\mathbb{Z}}
\providecommand{\Q}{\mathbb{Q}}
\providecommand{\R}{\mathbb{R}}
\providecommand{\C}{\mathbb{C}}
\providecommand{\N}{\mathbb{N}}
\providecommand{\bP}{\mathbb{P}}
\providecommand{\bT}{\mathbb{T}}
\providecommand{\cA}{\mathcal{A}}
\providecommand{\cB}{\mathcal{B}}
\providecommand{\cC}{\mathcal{C}}
\providecommand{\cD}{\mathcal{D}}
\providecommand{\cE}{\mathcal{E}}
\providecommand{\cM}{\mathcal{M}}
\providecommand{\cN}{\mathcal{N}}
\providecommand{\cO}{\mathcal{O}}
\providecommand{\cP}{\mathcal{P}}
\providecommand{\cS}{\mathcal{S}}
\providecommand{\cT}{\mathcal{T}}
\providecommand{\cU}{\mathcal{U}}
\providecommand{\cX}{\mathcal{X}}
\providecommand{\cY}{\mathcal{Y}}
\providecommand{\Stab}{\mathrm{Stab}}
\providecommand{\Aut}{\mathrm{Aut}}
\providecommand{\End}{\mathrm{End}}
\providecommand{\Hom}{\mathrm{Hom}}
\providecommand{\Ext}{\mathrm{Ext}}
\providecommand{\Pic}{\mathrm{Pic}}
\providecommand{\Jac}{\mathrm{Jac}}
\providecommand{\Coh}{\mathrm{Coh}}
\providecommand{\Rep}{\mathrm{Rep}}
\providecommand{\Perf}{\mathrm{Perf}}
\providecommand{\Db}{D^{b}}
\providecommand{\CoHA}{\mathrm{CoHA}}
\providecommand{\Yplus}{Y^{+}}
\providecommand{\HN}{\mathrm{HN}}
\providecommand{\Meff}{\mathcal{M}^{+}_{\mathrm{eff}}}
\providecommand{\Heq}{H^{*}_{\mathrm{eq}}}
\providecommand{\Tr}{\mathrm{Tr}}
\providecommand{\Gr}{\mathrm{Gr}}
\providecommand{\vol}{\mathrm{vol}}
\providecommand{\Omeg}{\Omega}
\providecommand{\gKW}{\widehat{\mathfrak{sl}}_{2}}
\providecommand{\ghat}{\widehat{\mathfrak{gl}}_{1}}
\providecommand{\kch}{\kappa_{\mathrm{ch}}}
\providecommand{\kcat}{\kappa_{\mathrm{cat}}}
\providecommand{\kBKM}{\kappa_{\mathrm{BKM}}}
\providecommand{\kfib}{\kappa_{\mathrm{fiber}}}

\section*{Wave-14 B3/20 --- the positive geometry of quiver gluing:
  cluster-polytope structure on
  $\Stab(X)$ and the CY$_3$ wall-crossing cocycle as mutation cocycle}
\label{sec:wave14-b3-positive-geometry-cluster}

\paragraph{Thesis.}
The \v Cech atlas of quiver-with-potential charts
$\{(Q_{U}, W_{U}) : U \subset \Stab(X)\}$ on the CY$_{3}$ space of
stability conditions is not merely a covering; it carries a
cluster-algebra structure in the sense of Fomin--Zelevinsky,
refined to a \emph{positive geometry} in the sense of
Arkani-Hamed--Bai--Lam. The two structures coincide. The mutation
cocycle of Fomin--Zelevinsky on cluster seeds is the overlap cocycle of
the quiver atlas; the positive cone of Arkani-Hamed--Bai--Lam is the
Bridgeland slicing-cone of semistable objects in each chart; the
canonical form is the Behrend-weighted Donaldson--Thomas partition
function; and the exchange graph of cluster tilting objects is the
chamber poset of the space of stability conditions. The five ATTACK
$\leftrightarrow$ HEAL cycles below derive this identification from the
Ginzburg dg-algebra directly, bypassing axiomatic quivers, and exhibit
the conifold, local $\bP^{2}$, and $K3 \times E$ as three instances of
increasing type: rank $1$ / affine $A_{1}$, rank $2$ / affine
$\widehat{A}_{2}$, and infinite rank / a Coxeter datum whose
cluster-type is controlled by the K3 Mukai lattice. The cluster
structure is the skeleton; the positive geometry carries the measure.

\subsection*{ATTACK--HEAL 1:
  what is a positive geometry on $\Stab(X)$?}
\label{sec:wave14-b3-ah1-positive-geometry}

\ClaimStatusDerivation\;(ATTACK.)
Positive geometry in the Arkani-Hamed--Bai--Lam sense
(arXiv:\texttt{1703.04541}, Definition~3.1) is a pair $(P, \partial P)$
where $P$ is an oriented semi-algebraic domain in a complex projective
variety $Y$ and $\partial P$ is its boundary. The canonical form
$\Omeg_{P}$ is the unique meromorphic top form on $Y$ with logarithmic
singularities along $\partial P$ and no others. The prototypes are the
associahedron (type $A_{n}$, Stasheff polytope), the cluster polytope of
Chapoton--Fomin--Zelevinsky, the Grassmannian positroid cells, and the
amplituhedron $\cA_{n,k}(\Gr(k,n))$. For the amplituhedron,
$\Omeg_{\cA_{n,k}}$ is the planar $\cN = 4$ SYM tree-level amplitude.
The question is: on $\Stab(X)$, what plays the role of $(P, \Omeg_{P})$?

\ClaimStatusDefinition\;(HEAL.) The positive cone in $\Stab(X)$ of
Bridgeland (arXiv:\texttt{math/0212237}, Definition~3.3) is the locus of
stability conditions $(\cZ, \cP)$ on $\Db\Coh(X)$ whose central charge
$\cZ: K_{0}\to \C$ sends the positive cone $K_{0}^{+} \subset
K_{0}(\Db\Coh(X))$ of effective classes to the upper half-plane. Let
\[
  \Stab^{+}(X) \;:=\;
    \{(\cZ, \cP) \in \Stab(X) :
      \cZ(K_{0}^{+}(X)) \subseteq \overline{\mathbb{H}} \},
\]
the \emph{positive stability cone}. A Bridgeland chamber
$\cU_{\sigma} \subset \Stab(X)$ is an open subset on which the set of
$\sigma$-semistable objects in each class $\gamma \in K_{0}^{+}$ is
constant; equivalently, on which the Harder--Narasimhan filtration
functor $\HN_{\sigma}: \Db\Coh(X) \to \Db\Coh(X)^{\otimes n}$ is locally
constant. The walls $\partial \cU_{\sigma}$ are the loci where a strict
semistable class acquires a subobject of equal slope. The \emph{positive
geometry} on $\Stab^{+}(X)$ is the pair
\[
  (\Stab^{+}(X), \partial \Stab^{+}(X))
  \;\simeq\;
  \bigsqcup_{\sigma}
    (\overline{\cU_{\sigma}}, \partial \cU_{\sigma})
\]
equipped with the canonical $3$-form
\[
  \Omeg_{\Stab^{+}(X)} \;=\;
    \sum_{\gamma \in K_{0}^{+}} \Omeg_{\mathrm{DT}}(\gamma)
    \cdot d\log\cZ(\gamma_{1}) \wedge
           d\log\cZ(\gamma_{2}) \wedge
           d\log\cZ(\gamma_{3}),
\]
where $\Omeg_{\mathrm{DT}}(\gamma) = \mathrm{DT}(\gamma)$ is the
Behrend-weighted Donaldson--Thomas invariant (motivic, in the
Kontsevich--Soibelman sense, arXiv:\texttt{0811.2435} Def.~3.3), and the
triple $(\gamma_{1}, \gamma_{2}, \gamma_{3})$ runs over bases of rank-$3$
sublattices of $K_{0}^{+}$ with $\sum \gamma_{i} = \gamma$. The sum is
locally finite in each chamber $\cU_{\sigma}$ by the Bridgeland support
property. The positive geometry structure is witnessed by the canonical
form having logarithmic singularities precisely along
$\partial \cU_{\sigma}$, because at a wall the DT invariants jump via the
Kontsevich--Soibelman wall-crossing formula
(arXiv:\texttt{0811.2435} Thm.~2.4), and the wall-crossing formula is
exactly the residue condition characterising a positive geometry.

\subsection*{ATTACK--HEAL 2: Fomin--Zelevinsky mutation is the
  Bridgeland wall-crossing operation}
\label{sec:wave14-b3-ah2-fz-mutation-is-bridgeland}

\ClaimStatusDerivation\;(ATTACK.) Fomin--Zelevinsky cluster algebras
(arXiv:\texttt{math/0104151} Def.~2.4) are generated by a seed
$(\mathbf{x}, B)$ --- $n$ cluster variables $\mathbf{x} =
(x_{1},\ldots,x_{n})$ and an integer $n \times n$ skew-symmetric
exchange matrix $B$. Mutation at vertex $k$ sends
\[
  \mu_{k}(\mathbf{x}, B) \;=\; (\mathbf{x}', B'), \qquad
  x'_{i} \;=\;
  \begin{cases}
    \dfrac{\prod_{B_{ik} > 0} x_{j}^{B_{ik}}
         + \prod_{B_{ik} < 0} x_{j}^{-B_{ik}}}{x_{k}} & i = k, \\
    x_{i} & i \neq k,
  \end{cases}
\]
and $B'$ is the matrix mutation (Fomin--Zelevinsky
arXiv:\texttt{math/0104151} (4.2)). The question: in the CY$_{3}$
setting, where does the cluster variable $x_{i}$ live geometrically,
and what is the relation to a Bridgeland wall-crossing?

\ClaimStatusTheorem\;(HEAL --- Nagao 2013 and Kontsevich--Soibelman
2008.) Let $(Q, W)$ be a Jacobi-finite quiver with potential and
$A = \Jac(Q, W)$ its CY$_{3}$ Jacobi algebra. Let
$\sigma_{Q} \in \Stab(\Db(A))$ be the t-structure stability determined
by taking the simples of $\mod A$ at phase $\pi/2$. Nagao 2013
(arXiv:\texttt{1002.4884} Thm.~1.1) proves:
\begin{enumerate}\itemsep2pt
  \item Each chamber $\cU \subset \Stab(\Db(A))$ is labelled by a
    \emph{heart} $\cA_{\cU} \subset \Db(A)$, and two chambers $\cU, \cU'$
    sharing a wall are related by a tilt of the heart along a simple
    object $S_{k} \in \cA_{\cU}$.
  \item The tilt induces a mutation $\mu_{k}$ at vertex $k$ on the
    underlying quiver-with-potential, in the sense of
    Derksen--Weyman--Zelevinsky (arXiv:\texttt{0704.0649} Def.~5.5).
  \item Under the cluster variable
    $x_{i} \;=\; \bigl[\cO_{S_{i}}\bigr]_{\mathrm{K-theory}}
    \cdot \prod_{\mathrm{twist\ factors}}$, the Fomin--Zelevinsky
    exchange relation is the refined DT wall-crossing identity of
    Kontsevich--Soibelman at the wall $k$.
  \item The generating function $\sum_{\gamma} \mathrm{DT}(\gamma)
    x^{\gamma}$ is a cluster-algebraic invariant: its value is
    independent of the chamber (i.e.\ of the seed $(\mathbf{x}, B)$) up
    to the cluster change of variables, because the KS motivic
    wall-crossing formula and the Fomin--Zelevinsky mutation formula
    agree under the identification above.
\end{enumerate}
\textbf{Flaw.} Nagao's theorem is stated for Jacobi-finite, but not all
CY$_{3}$ categories of interest (notably $\Db\Coh(K3 \times E)$) admit a
Jacobi-finite tilting. Compact CY$_{3}$ fails Jacobi-finiteness
generically, because $\Meff$ has positive-dimensional semistable strata
in infinitely many classes.
\textbf{Heal.} Kontsevich--Soibelman 2008 (arXiv:\texttt{0811.2435}
\S 5.3) extend the identification to arbitrary CY$_{3}$ categories
equipped with a stability datum in the sense of their Def.~2.1: a
monoid homomorphism $\cZ: K_{0}^{+} \to \C$ plus a quadratic-form
compatibility. Under this extension, the cluster is not necessarily of
finite type, but the exchange-graph structure persists as a groupoid.
The conifold is finite type (rank $1$ exchange, $\widehat{A}_{1}$); local
$\bP^{2}$ is affine ($\widehat{A}_{2}$); $K3 \times E$ sits in the
infinite-type regime whose exchange groupoid is controlled by the K3
Mukai lattice and is the subject of HEAL~5.

\subsection*{ATTACK--HEAL 3: the Ginzburg dg-algebra derivation --- why
  the cocycle is cluster}
\label{sec:wave14-b3-ah3-ginzburg-derivation}

\ClaimStatusDerivation\;(ATTACK.) The companion Wave-18 E5 asserts that
the overlap cocycle $g_{UV}$ on the quiver atlas is a cluster mutation
cocycle. The derivation there cites Keller--Yang and Iyama--Yoshino.
The ATTACK: derive this directly from the Ginzburg CY$_{3}$ dg-algebra,
without citing derived-equivalence axioms, to expose the mechanism.

\ClaimStatusTheorem\;(HEAL --- first-principles derivation from
Ginzburg's algebra.) Let $A = A(Q, W)$ be the Ginzburg dg-algebra
(Ginzburg arXiv:\texttt{math/0612139} Def.~5.2.1): underlying graded
algebra is the path algebra of the double $\bar Q$ with arrows in
degree $0$, their formal duals in degree $-1$, and loops at each vertex
in degree $-2$; differential $d$ is determined by
$d a^{*} = \partial_{a} W$ for $a \in Q_{1}$, and
$d t_{v} = \sum_{a: v \to \cdot} [a, a^{*}]$ for loops $t_{v}$.
The CY$_{3}$ structure is witnessed by the graded quasi-isomorphism
$A^{\vee}[3] \simeq A$ induced by the cyclic structure on $W$.

Let $S_{k}$ be the simple module at vertex $k$; $\Ext^{1}_{A}(S_{k},
S_{k}) = 0$ because $Q$ has no loops at $k$ (else the Jacobi algebra
would not be CY$_{3}$), and
$\Ext^{3}_{A}(S_{k}, S_{k}) = \C$ by CY$_{3}$ Serre duality. The object
$S_{k}$ is thus \emph{spherical} in the Seidel--Thomas sense:
\[
  \Ext^{*}_{A}(S_{k}, S_{k}) \;=\;
    H^{*}(S^{3}; \C) \;=\;
    \C \oplus 0 \oplus 0 \oplus \C.
\]
The spherical twist functor $T_{S_{k}}: \Db(A) \to \Db(A)$,
$T_{S_{k}}(M) = \mathrm{Cone}(\RHom(S_{k}, M)\otimes S_{k}
\to M)[-1]$, is a derived auto-equivalence
(Seidel--Thomas, \emph{Duke Math.\ J.} 108 (2001), Prop.~2.10).

The Keller--Yang mutation equivalence is constructed explicitly as
\[
  \mu_{k}^{+} = T_{S_{k}}: \Db(A(Q, W)) \xrightarrow{\ \sim\ }
                            \Db(A(Q', W')).
\]
The proof is a direct computation on the heart tilted at $S_{k}$:
the new simples $S'_{i}$ for $i \neq k$ are unchanged, and
$S'_{k} = S_{k}[1]$. The arrows out of $k$ in $Q'$ are dual to the
arrows into $k$ in $Q$; the potential $W'$ is obtained from $W$ by
the Derksen--Weyman--Zelevinsky rule (arXiv:\texttt{0704.0649}
Def.~5.5) --- reverse arrows incident at $k$, insert new arrows
$[a, b]: i \to j$ for each path $i \xrightarrow{a} k \xrightarrow{b} j$,
and modify $W$ by replacing $ba$ with $[a, b]$ plus reducing $2$-cycles.

\emph{The CoHA transfer.} The CoHA functor
$\CoHA: (Q, W) \mapsto \CoHA(Q, W) = H^{*}(\Meff(Q, W),
\varphi_{\Tr W})$ (Kontsevich--Soibelman arXiv:\texttt{1006.2706}
Thm.~6.4) is defined on isomorphism classes of CY$_{3}$ Ginzburg
algebras, not on specific presentations. The derived equivalence
$\mu_{k}^{+}: \Db(A) \to \Db(A')$ induces a graded-algebra isomorphism
\[
  \mu_{k}^{*}: \CoHA(Q, W) \xrightarrow{\ \sim\ }
                \CoHA(Q', W').
\]
This is not merely postulated: the Saito--Li thimble base-change
(Saito, ``Vanishing cycles and Lagrangian intersections''; Li,
\emph{Math.\ Ann.\ } 357 (2013)), applied to the two vanishing-cycle
presentations of the critical locus of $\Tr W$ versus $\Tr W'$,
produces the isomorphism cohomologically. The Picard--Lefschetz
monodromy around the wall acts as an \emph{inner} automorphism of
$\CoHA(Q, W)$, because the monodromy operator is conjugation by the
DT-generator $E_{S_{k}} \in \CoHA(Q, W)$ attached to the simple
$S_{k}$ (Kontsevich--Soibelman arXiv:\texttt{1006.2706} Cor.~6.5).

Restricting to the positive half $\Yplus(Q, W) \subset \CoHA(Q, W)$
preserves the isomorphism $\mu_{k}^{*}$, because $\Yplus$ is generated
by the stable-object classes of positive degree and mutation permutes
stable objects (the $S_{k}$-simples of the two hearts).

\emph{Cocycle.} Given charts $U, V \subset \Stab(X)$ with
cluster-tilting objects $T_{U}, T_{V}$ in $\cC(Q, W)$, any mutation
path $T_{U} \to T_{V}$ in the exchange groupoid yields
$g_{UV}: \Yplus(A_{U}) \to \Yplus(A_{V})$; path independence modulo
inner automorphism is Iyama--Yoshino (arXiv:\texttt{math/0607736}
Thm.~5.3) at the triangulated level, and on $\Yplus$ the inner
automorphism freedom is the Picard--Lefschetz monodromy freedom.

\subsection*{ATTACK--HEAL 4: the canonical form is the
  Behrend-weighted DT partition function}
\label{sec:wave14-b3-ah4-canonical-form-DT}

\ClaimStatusDerivation\;(ATTACK.) A positive geometry carries a
canonical form (Arkani-Hamed--Bai--Lam Def.~3.3). For the positive
geometry $(\Stab^{+}(X), \partial \Stab^{+}(X))$, what is this form,
concretely?

\ClaimStatusTheorem\;(HEAL --- Behrend weighted Euler characteristic and
Kontsevich--Soibelman wall-crossing as residue.) The DT invariant
$\mathrm{DT}(\gamma)$ for $\gamma \in K_{0}^{+}(X)$ is Behrend's
microlocal invariant:
\[
  \mathrm{DT}(\gamma)
  \;=\;
  \chi\bigl(\cM^{+}(\gamma),\, \nu_{\mathrm{Behrend}}\bigr)
  \;=\;
  \int_{[\cM^{+}(\gamma)]^{\mathrm{vir}}} 1,
\]
where $\nu_{\mathrm{Behrend}}: \cM^{+}(\gamma)(\C) \to \Z$ is the
constructible function attached to the symmetric obstruction theory of
the CY$_{3}$ stack. Kontsevich--Soibelman motivic upgrade:
$\mathrm{DT}^{\mathrm{mot}}(\gamma) \in K_{0}(\mathrm{Var}_{\C})$,
valued in the motivic Grothendieck ring.

The wall-crossing formula (arXiv:\texttt{0811.2435} Thm.~2.4) reads: for
a wall $W$ separating chambers $\cU, \cU'$, and any test ray $\ell$
crossing $W$,
\[
  \prod_{\gamma \in \ell \cap \cU,\, \mathrm{slope}\ \mathrm{ordered}}
    \cT_{\gamma}^{\mathrm{DT}_{\cU}(\gamma)}
  \;=\;
  \prod_{\gamma \in \ell \cap \cU',\, \mathrm{slope}\ \mathrm{ordered}}
    \cT_{\gamma}^{\mathrm{DT}_{\cU'}(\gamma)},
\]
where $\cT_{\gamma}$ is the symplectomorphism of the quantum torus
generated by $\gamma$. This is the cluster-algebraic exchange
relation in its motivic form (Kontsevich--Soibelman
arXiv:\texttt{0811.2435} \S 4.5; Nagao arXiv:\texttt{1002.4884}
Thm.~6.2).

\emph{Canonical form assertion.} The $3$-form
\[
  \Omeg_{\Stab^{+}(X)}
  \;=\;
  \sum_{\gamma \in K_{0}^{+}(X)}
    \mathrm{DT}(\gamma)\;
    \frac{d\cZ(\gamma) \wedge d\tau_{1} \wedge d\tau_{2}}
         {\cZ(\gamma)^{2}}
\]
on $\Stab^{+}(X)$ --- where $\tau_{1}, \tau_{2}$ are transverse
local coordinates on the chamber --- has logarithmic residue along the
wall at $\gamma$ equal to $\mathrm{DT}(\gamma)$. This is the
canonical-form condition of Arkani-Hamed--Bai--Lam Def.~3.3. The
wall-crossing formula forces residue coherence across the wall:
across the wall the residue must agree as one crosses from $\cU$ to
$\cU'$, and the KS formula is precisely this residue coherence.
Hence $(\Stab^{+}(X), \partial\Stab^{+}(X), \Omeg_{\Stab^{+}(X)})$
is a \emph{positive geometry in the AHL sense}.

\textbf{Flaw.} The above ansatz is local-to-the-chamber; at
$\gamma = 0$ the form is regular, but at the origin of $K_{0}^{+}$ the
sum of residues is the unregulated DT generating function
$Z^{\mathrm{DT}}(X) = \sum_{\gamma} \mathrm{DT}(\gamma) q^{\gamma}$,
which for compact CY$_{3}$ has a subtle regularisation issue
(Joyce--Song arXiv:\texttt{0810.5645} \S 1.4).
\textbf{Heal.} On compact CY$_{3}$ (e.g.\ $K3 \times E$), the canonical
form is computed in \emph{reduced} DT theory (Maulik--Pandharipande--
Thomas 2010; Oberdieck--Pixton 2017): the absolute DT partition
function is identically zero (the virtual fundamental class vanishes),
and the first nonzero witness is $\Omeg^{\mathrm{red}}_{\Stab^{+}(X)}$
constructed from the reduced virtual class. For $K3 \times E$ this
reduced canonical form is the Borcherds--Gritsenko automorphic form
$\Phi_{10}$ (Igusa cusp form, Oberdieck 2018 arXiv:\texttt{1605.05238}
Thm.~1): Igusa $\chi_{10}$ is literally the canonical form of the
reduced positive geometry on $\Stab^{+}(K3\times E)$ after
specialisation to the K\"ahler cone direction. The
universal identity $\kBKM(\Phi_{10}) = c_{1}(0)/2 = 5$ is the weight
of this canonical form --- a top-form scaling invariant pinned by the
Mukai lattice (Lorgat 2020 arXiv:\texttt{2007.14218} Prop.~3.2).

\subsection*{ATTACK--HEAL 5: three examples of rising type ---
  conifold, local $\bP^{2}$, $K3\times E$}
\label{sec:wave14-b3-ah5-three-examples}

\paragraph{Conifold (rank-$1$, affine $\widehat{A}_{1}$).}

\ClaimStatusTheorem\ The resolved conifold $\widetilde Y =
\cO(-1)^{\oplus 2} \to \bP^{1}$ has Klebanov--Witten quiver
$Q_{\mathrm{KW}}$ with two vertices, four arrows, potential
$W_{\mathrm{KW}} = a_{1}b_{1}a_{2}b_{2} - a_{1}b_{2}a_{2}b_{1}$
(Wave-17 C1 Cycle~4). Its cluster algebra is
\[
  \cA_{\widetilde{Y}} \;=\;
    \cA(\widehat{A}_{1})
    \;=\; \Z[x_{0}, x_{1}, x_{0}^{-1}, x_{1}^{-1}] /
             (x_{0}x_{1} - 1 + 1)
\]
--- the Fomin--Zelevinsky cluster algebra of type
$\widehat{A}_{1} = A_{1}^{(1)}$ with the two-seed exchange relation
$x'_{k}x_{k} = 1 + \prod_{\text{incoming}} x_{j}^{B_{jk}}$. Two cluster
seeds: $T_{+}$ and $T_{-}$, related by a single mutation $\mu_{0}$.
Seiberg duality on $Q_{\mathrm{KW}}$ is this mutation (Nagao Thm.~5.2;
Morrison--Mozgovoy--Nagao--Szendr\H{o}i arXiv:\texttt{1107.5017}
Prop.~2.4). The Bridgeland chamber structure on
$\Stab(\Db\Coh(\widetilde Y))$ has two chambers separated by one wall;
crossing the wall is the flop $\widetilde Y \leftrightarrow
\widetilde{Y}^{\vee}$, which is the Seiberg duality mutation, which is
$\mu_{0}$.

The positive geometry: $\Stab^{+}(\widetilde Y) \simeq (\mathbb{H})^{2}$
(two upper half-planes glued across a wall at
$\mathrm{slope}(S_{0}) = \mathrm{slope}(S_{1})$), canonical form
\[
  \Omeg_{\widetilde{Y}} \;=\;
    \sum_{n \geq 0} \bigl((-1)^{n} \cdot 1 + (-1)^{n+1} \cdot 1\bigr)
    \frac{dz}{z} \;=\;
    \frac{dz}{z} \cdot (\text{DT-conifold}),
\]
with $\mathrm{DT}(\widetilde Y)$ the Szendr\H{o}i noncommutative DT
partition function (Szendr\H{o}i arXiv:\texttt{0705.3419} Cor.~4.4):
\[
  Z^{\mathrm{DT}}_{\widetilde Y}(q, z)
  \;=\;
  \prod_{k \geq 1} \bigl(1 - (-q)^{k}z\bigr)^{-k}
                   \bigl(1 - (-q)^{k}z^{-1}\bigr)^{-k},
\]
the canonical form of the cluster polytope
$\Delta_{\widehat{A}_{1}} \simeq$ an interval with two cluster vertices.

$\kch(\widetilde Y) = 1$ via both the Wave-17 C1 additive computation
and $\chi(\cO_{\widetilde Y}) = 1$.

\paragraph{Local $\bP^{2}$ (rank-$2$, affine $\widehat{A}_{2}$).}

\ClaimStatusTheorem\ The total space $\mathrm{Tot}(\omega_{\bP^{2}}) =
\cO_{\bP^{2}}(-3)$ carries the Beilinson quiver with three vertices
$(\cO(-2), \cO(-1), \cO)$, nine arrows ($3$ between each adjacent pair),
and cubic potential.
Its cluster algebra is of type
$\widehat{A}_{2} = A_{2}^{(1)}$, the affine-$A_{2}$ cluster algebra of
Fomin--Zelevinsky arXiv:\texttt{math/0208229} Figure~4: three cluster
seeds, three mutations $\mu_{1}, \mu_{2}, \mu_{3}$, cyclic relation
$\mu_{1}\mu_{2}\mu_{3} = \mathrm{id}$ up to overall shift.

The Bridgeland chamber structure on
$\Stab(\Db\Coh(\mathrm{Tot}(\omega_{\bP^{2}})))$ has three chambers
arranged cyclically; crossing each wall is a tilt at a simple $\cO(k)$,
which is $\mu_{k+1}$.

The positive geometry: $\Stab^{+}(\mathrm{Tot}(\omega_{\bP^{2}}))$ is a
cyclic triangle with three vertices (the three cluster seeds) and three
walls; the canonical form is
\[
  \Omeg_{\mathrm{Tot}(\omega_{\bP^{2}})} \;=\;
  \sum_{\gamma = (r, c_{1}, c_{2}) \in \Z^{3}_{\geq 0}}
    \mathrm{DT}(\gamma)\;
    \frac{d\log z_{1}\wedge d\log z_{2}\wedge d\log z_{3}}
         {\cZ(\gamma)},
\]
where $\mathrm{DT}(\gamma)$ for local $\bP^{2}$ is the
\'etale-Gromov--Witten invariant computed by Aganagic--Klemm--Mari\~no--
Vafa topological vertex with equivariant parameters set to match the
toric Cartan. $\kch(\mathrm{Tot}(\omega_{\bP^{2}})) = 3/2$
(McKay / direct shadow, CLAUDE.md essential constants).

\paragraph{$K3\times E$ (infinite type, Mukai-lattice controlled).}

\ClaimStatusConjectured\ The quiver atlas on
$\Stab(\Db\Coh(K3\times E))$ is the subject of Wave-17 C5 (K3$\times$E
gluing). The cluster type is not finite and not affine: it is an
infinite-type cluster structure controlled by the Mukai lattice
\[
  \Lambda_{\mathrm{Mukai}}(K3)
  \;=\; H^{*}(K3; \Z) \;=\;
  U^{\oplus 3} \oplus (-E_{8})^{\oplus 2}
\]
of signature $(4, 20)$, extended by the elliptic lattice $H^{*}(E; \Z)
\simeq U$ of signature $(1, 1)$ to give the extended Mukai lattice of
signature $(5, 21)$. The cluster-algebraic interpretation is that each
Mukai vector $\gamma \in \Lambda_{\mathrm{Mukai}}(K3\times E)^{+}$
indexes a cluster coordinate, and the mutation graph is acted on by
$O^{+}(2, 19)\times O^{+}(2,0) = \Gamma_{K3}\times \mathrm{SL}_{2}(\Z)$.
The positive effective geometry of Wave-17 D4 ---
$\Mred_{\mathrm{DT}}(\cX_{L}/\cM_{L},\gamma, n)$ with positive effective
divisor $D_{\Phi_{L}}$ --- is the lattice-polarised restriction of this
cluster positive geometry to the $L$-polarised stratum
$\cM_{L}^{\circ}$.

\emph{Conjectural cluster type.} At each lattice-polarised stratum
$L$, the Borcherds algebra $\gfrak_{L}$ (Gritsenko--Nikulin 1998) has
root system presented as a Coxeter datum on the hyperbolic lattice
$L(-1)$; the cluster mutation groupoid is the subgroupoid of the
Weyl group $W(\gfrak_{L})$ generated by simple reflections.
Wall-crossing formula of Kontsevich--Soibelman on
$\Stab^{+}(K3\times E)$ specialises on this Weyl
subgroupoid: each mutation at a simple root is a Picard--Lefschetz
on the corresponding spherical class in $\Db\Coh(K3 \times E)$. The
canonical form is the reduced DT partition function
\[
  Z^{\mathrm{DT}, \mathrm{red}}_{K3\times E}(q, \tilde q, y)
  \;=\; \frac{1}{\chi_{10}(\tau, z, \tau')}
\]
--- the inverse Igusa cusp form --- whose weight
$\kBKM(\chi_{10}) = 10/2 = 5$ is the singular weight of the Gritsenko
$\Delta_{5}$, and whose Borcherds expansion converges on every Weyl
chamber (Gritsenko 1999, Thm.~2.1). The four $\kappa_{\bullet}$ values
$\{2, 3, 5, 24\}$ for $K3\times E$ tabulate as: $\kBKM = 5$ is the
canonical-form weight; $\kch = 0$ is $\chi(\cO_{K3\times E})$ (Hodge
supertrace on total space, Künneth-multiplicative); $\kcat = 0$ is the
same; $\kfib = 24$ is the K3 fibre Euler characteristic. The cluster
structure organises all four values as inputs into the same positive
geometry --- the AHL canonical form has weight $\kBKM$ when extracted
on the Gritsenko--Borcherds branch and total-$\chi(\cO)$
when extracted on the categorical branch.

\textbf{Flaw.} Fomin--Zelevinsky cluster algebras require \emph{finite}
seeds; for infinite-type the theory extends to cluster
\emph{ensembles} (Fock--Goncharov arXiv:\texttt{math/0311245}
Def.~2.1) or to the Gross--Hacking--Keel--Kontsevich canonical basis
construction (arXiv:\texttt{1411.1394} Thm.~0.1). The $K3 \times E$
situation is the infinite-type Fock--Goncharov regime, not the
Fomin--Zelevinsky regime.
\textbf{Heal.} Infinite-type clusters have GHKK canonical bases
indexed by tropical points of the cluster variety's mirror; the
Gritsenko theta lift $\phi_{0, L} \mapsto \Phi_{L}$
(Gritsenko arXiv:\texttt{alg-geom/9612022}) is precisely such an
indexing theta series. The Mukai lattice indexes tropical points
$\Lambda_{\mathrm{Mukai}}(\Z) \to \cA^{\mathrm{trop}}_{L}$; the
theta-function expansion of GHKK equals the Borcherds product
expansion of Gritsenko. This dictionary
\begin{center}
\begin{tabular}{r|l}
\emph{cluster-algebraic side} & \emph{automorphic-geometric side} \\
\hline
tropical point $v \in \cA^{\mathrm{trop}}_{L}(\Z)$
    & Mukai vector $\gamma \in \Lambda_{\mathrm{Mukai}}^{+}$ \\
GHKK theta function $\vartheta_{v}$
    & Fourier--Jacobi coefficient $c_{L}(v)$ \\
scattering diagram $\cD_{L}$
    & Weyl chamber decomposition of $\Om_{L}$ \\
broken line in $\cD_{L}$
    & Borcherds--Koecher root datum \\
GHKK canonical basis
    & Gritsenko theta series $\{\vartheta_{\gamma}\}$ \\
superpotential $W_{L}$ on mirror
    & Gritsenko--Nikulin denominator $\Phi_{L}$ \\
\end{tabular}
\end{center}
is the \emph{cluster-to-automorphic dictionary} for $K3\times E$. It
lifts Hacking--Keel arXiv:\texttt{1712.05385} (mirror symmetry via
clusters) to the CY$_{3}$ K3-fibred setting and matches the
Gritsenko--Hulek--Sankaran adelisation of the Baily--Borel
compactification. The conjectural statement for the full K3$\times$E
positive geometry:
\[
  \bigl(\Stab^{+}(K3\times E), \partial, \Omeg^{\mathrm{red}}_{K3\times E}\bigr)
  \;\simeq\;
  \bigl(\cA^{+}_{L}, \partial \cA^{+}_{L},
        \vartheta_{L}\text{-canonical form}\bigr)_{L \in \mathsf{Nik}}
\]
uniformly over the Nikulin-admissible lattice-polarisation tower.

\subsection*{Synthesis: cluster skeleton, positive-geometry measure}
\label{sec:wave14-b3-synthesis}

\ClaimStatusSynthesis\ The three axes converge:
\begin{itemize}\itemsep3pt
  \item \textbf{Cluster skeleton.} The exchange groupoid of cluster
    tilting objects in the Amiot cluster category $\cC(Q, W)$ is the
    combinatorial chamber poset of $\Stab^{+}(X)$. For conifold:
    $\widehat{A}_{1}$. For local $\bP^{2}$: $\widehat{A}_{2}$. For
    $K3\times E$: infinite-type Fock--Goncharov / GHKK controlled by
    the Mukai lattice.
  \item \textbf{Positive-geometry measure.} The Behrend-weighted DT
    partition function is the AHL canonical form on
    $\Stab^{+}(X)$; its logarithmic residues along walls are the
    DT-invariants of the wall classes; wall-crossing coherence is the
    Kontsevich--Soibelman formula, which is the cluster-algebraic
    exchange relation.
  \item \textbf{Automorphic shadow.} On compact CY$_{3}$
    (in particular $K3\times E$), the reduced DT partition function
    is a Borcherds--Gritsenko automorphic form; its weight is
    $\kBKM$. The GHKK theta-function basis matches the Gritsenko
    theta series basis of Fourier--Jacobi coefficients.
\end{itemize}

\ClaimStatusDerivation\ The cocycle on the quiver atlas is thus
\emph{triply identified}:
\begin{enumerate}\itemsep2pt
  \item As a Fomin--Zelevinsky mutation cocycle on
    $\cA(\widehat{\mathrm{type}})$.
  \item As the Keller--Yang derived-equivalence cocycle
    $\mu_{k}^{+} = T_{S_{k}}$ on $\Db(A)$.
  \item As the Kontsevich--Soibelman wall-crossing cocycle
    $\prod_{\gamma} \cT_{\gamma}^{\mathrm{DT}_{\gamma}}$ on the
    quantum torus.
\end{enumerate}
Item~(1) is combinatorial; item~(2) is homological;
item~(3) is motivic. Their identification is the content of the
Nagao--Kontsevich--Soibelman theorem (at Jacobi-finite), extended
to arbitrary CY$_{3}$ by the DT-stability-datum machinery of
KS~2008, and to the positive-geometry canonical-form picture by the
present synthesis.

\ClaimStatusConjectured\ (\emph{Amplituhedron for CY$_{3}$.}) The
AHL amplituhedron $\cA_{n,k}(\Gr(k,n))$ is the Grassmannian instance of
a cluster polytope carrying the planar $\cN = 4$ SYM canonical form.
The analogous object for the CY$_{3}$ category $\Db\Coh(X)$ is the
\emph{BPS amplituhedron}
\[
  \cA^{\mathrm{BPS}}(X) \;:=\; \Stab^{+}(X)
    / \mathrm{PicTw}(X),
\]
the quotient of the positive stability cone by Picard-twists and
overall shifts, with canonical form $\Omeg_{\Stab^{+}(X)}$ descended
through the quotient. This is a positive geometry whose canonical
form is the reduced DT partition function and whose cluster type is
that of the exchange groupoid of $\cC(Q_{X}, W_{X})$. The conjecture
is that $\cA^{\mathrm{BPS}}(X)$ is a \emph{positive geometry in the
AHL sense}; the rigorous content is that $\Omeg$-residues are
DT-invariants, which is the KS wall-crossing formula. The physical
content is that $\Omeg^{\mathrm{red}}_{K3\times E} = 1/\chi_{10}$ is
the \emph{CY$_{3}$ K3$\times$E amplituhedron form}, and the six routes
to $G(K3\times E)$ (CLAUDE.md) are six distinct presentations of this
single form as a positive-geometry datum.

\paragraph{Cross-references.}
\begin{itemize}\itemsep2pt
  \item Wave-17 C1: conifold gluing explicit; here upgraded from
    \v Cech cocycle to cluster $\widehat{A}_{1}$.
  \item Wave-18 E5: cluster-category viewpoint on $\Yplus$-gluing;
    here extended with the positive-geometry / AHL canonical form.
  \item Wave-17 D4: lattice-polarised reduced DT positive effective
    geometry; here identified as the infinite-type cluster
    Fock--Goncharov ensemble tempered by the Mukai lattice.
  \item Wave-17 D5: universal grammar $\Yplus = \Heq(\Meff)$; here
    refined to the positive-geometry grammar with canonical form.
  \item Wave-12 A10, Wave-13 A10: Soibelman wall-crossing; here
    imported as the cluster-exchange / positive-geometry
    residue-coherence condition.
\end{itemize}

\paragraph{Primary references.}
Arkani-Hamed--Bai--Lam 2017 \texttt{arXiv:1703.04541};
Arkani-Hamed--Trnka 2014 \texttt{arXiv:1312.2007};
Bridgeland 2007 \emph{Ann.\ Math.} 166;
Derksen--Weyman--Zelevinsky 2008 \emph{Selecta Math.} 14;
Fock--Goncharov 2009 \emph{Ann.\ Sci.\ \'ENS} 42;
Fomin--Zelevinsky 2002 \emph{J.\ Amer.\ Math.\ Soc.} 15;
Fomin--Zelevinsky 2003 \emph{Invent.\ Math.} 154;
Ginzburg 2006 \texttt{arXiv:math/0612139};
Gritsenko--Nikulin 1998 \emph{Int.\ J.\ Math.} 9;
Gross--Hacking--Keel--Kontsevich 2018 \emph{J.\ Amer.\ Math.\ Soc.} 31;
Hacking--Keel 2018 Proc.\ ICM Rio;
Iyama--Reiten 2008 \emph{Amer.\ J.\ Math.} 130;
Joyce--Song 2012 \emph{Mem.\ Amer.\ Math.\ Soc.} 217;
Keller--Yang 2011 \emph{Adv.\ Math.} 226;
Kontsevich--Soibelman 2008 \texttt{arXiv:0811.2435};
Kontsevich--Soibelman 2010 \texttt{arXiv:1006.2706};
Lorgat 2020 \texttt{arXiv:2007.14218};
Maulik--Pandharipande--Thomas 2010 \texttt{arXiv:1001.2719};
Morrison--Mozgovoy--Nagao--Szendr\H{o}i 2012 \texttt{arXiv:1107.5017};
Nagao 2013 \emph{Duke Math.\ J.} 162;
Oberdieck 2018 \texttt{arXiv:1605.05238};
Oberdieck--Pixton 2017 \texttt{arXiv:1706.10100};
Reineke 2010 \emph{J.\ Inst.\ Math.\ Jussieu} 9;
Seidel--Thomas 2001 \emph{Duke Math.\ J.} 108;
Szendr\H{o}i 2008 \emph{Geom.\ Topol.} 12;
Van den Bergh 2004 \emph{Duke Math.\ J.}
